\begin{center}
\textsc{Gauss} an \textsc{Bessel}.\quad Göttingen, 9.~April 1830.
\end{center}

\par\noindent\hspace*{\parindent}.\ .\ .\ .\ .\ .\ Seit Mai 1829 haben mir die trigonometrischen Messungen noch viel mehr Zeit geraubt, obwohl ich diesmal keinen unmittelbaren Anteil an den Arbeiten im Feld genommen habe.
Noch in dieser Stunde bin ich mit der Verarbeitung der vorjährigen Messungen nicht ganz fertig (Gott sei Dank aber beinahe), wobei ich ohne jede Hilfe auskommen muss.
Doch hat es mir viel Freude gemacht, dass mein Sohn seinen Anteil an diesen Arbeiten zu meiner größten Zufriedenheit ausgeführt hat.
Er hat ganz allein ein großes Dreiecksnetz von der Weser bis zur niederländischen Grenze geführt (über das Osnabrücker Gebiet, wobei er auch mehrfach preußisches Gebiet berührte und von den dortigen Behörden sehr großzügig unterstützt wurde), wodurch außer den Hauptpunkten noch etwa 250 weitere Punkte festgelegt sind.
Die --- von einem anderen Offizier --- im Fürstentum Hildesheim vorgenommenen trigonometrischen Messungen sind ebenfalls so gut wie abgeschlossen, sodass alle Messtischblätter (53, wovon etwa ein Viertel bereits aufgenommen ist) mit Festpunkten versehen werden können.\ .\ .\ .\ .\ .\ .

\begin{center}\rule{.24\linewidth}{.4pt}\end{center}
% The printer's signature is visible in the Google witness, physical leaf 372.
\par\smallskip\noindent\hfill{\footnotesize 46\textsuperscript{*}}\par
